% =============================================================
% Agents
% Chapter order is set in ../main.tex; the rendered chapter number
% comes from \chapter{} below.
%
% This is the entry point to the agents arc. Subsequent chapters
% on orchestration, multi-agent systems, and context engineering
% build on the vocabulary established here.
% =============================================================

\chapter{Agents}
\label{ch:agents}

\todo{Chapter scaffold. The first section (``What is an Agent?'')
is written; the rest are stubs indicating intended content. The
agent ecosystem survey in \S\ref{sec:agents-landscape} draws on
the AI Agent Index~\citep{staufer2026aiagentindex} and is one of
the load-bearing pieces of evidence for the book's design
recommendations.}

The previous chapters built up a working picture of the language
model as a stateless conditional sampler
(Chapter~\ref{ch:mental}), of the tools that extend its reach
beyond text generation (Chapter~\ref{ch:toolcalls}), and of the
retrieval machinery that supplies it with relevant context
(Chapter~\ref{ch:context}). What we have not yet done is assemble
those pieces into something that acts. This chapter does that
assembly. It opens with the terminological question --- what
exactly counts as an ``agent,'' given that the word is used
differently in every field that has ever borrowed it --- and
proceeds to the architectural one: how the model, the tool layer,
the memory subsystem, and the control loop combine into a system
that pursues goals over time.

\section{What is an Agent?}
\label{sec:agents-what}

The honest answer to the section title is that there is no single
definition, and the reader who has been around the AI literature
long enough will already have noticed that the term shifts shape
depending on who is using it. A control theorist, a robotics
researcher, a software engineer building microservices, an
economist modeling rational choice, a reinforcement-learning
researcher, and a philosopher of action all use the word
``agent,'' and they do not mean the same thing. Recent surveys of
the agentic-AI ecosystem are unusually candid about this: the
2025 AI Agent Index opens its background section by stating
flatly that ``definitions of AI agents are nebulous and differ
across fields,'' and traces the disagreement across at least six
disciplines~\citep{staufer2026aiagentindex}. That candor is a
useful starting point. It is more productive to acknowledge that
the term is contested than to legislate a definition the rest of
the field will ignore.

The disagreement is not arbitrary, however. Each field's usage
captures something real, and the differences are mostly about
which aspect of agency the field cares to foreground. A short
tour of the most consequential traditions, in roughly
chronological order, makes the picture legible.

\paragraph{Cybernetics.}
Norbert Wiener's mid-twentieth-century program took the
\emph{feedback loop} as the elementary unit of purposeful
behavior. A thermostat senses room temperature, compares it to a
setpoint, and acts on a heater to close the gap; a homing missile
senses bearing error and steers to drive it to zero; an organism
senses hunger and acts to reduce it. The cybernetic tradition
treats agency as the regulatory closure of such a loop:
\emph{sense, compare, act, repeat}, with the loop's goal encoded
in whatever the system is trying to drive toward zero. Nothing in
this account requires intelligence, language, or symbols. A
bimetallic strip in a thermostat counts as an agent under the
strict reading, which is both the tradition's strength
(everything purposeful fits one framework) and its embarrassment
(the framework includes things no one would intuitively call an
agent).

\paragraph{Artificial life and situated cognition.}
In the 1980s and 1990s, researchers including Christopher
Langton, Rodney Brooks, and Randall Beer pushed back against the
symbol-manipulation models that had dominated classical AI and
argued that agency emerges from \emph{situatedness}: a body, an
environment, and the closed sensorimotor loop between them.
Brooks's subsumption architecture built mobile robots whose
``intelligent'' behavior emerged from layered reactive
controllers with no internal world model. The slogan ``the world
is its own best model'' captures the spirit. Agency, in this
view, is not a property of a disembodied reasoner but of a
system embedded in and coupled to an environment that pushes back.

\paragraph{Rational agency.}
The textbook treatment most engineers first encounter is the one
in Russell and Norvig: an agent is anything that \emph{perceives}
its environment through sensors and \emph{acts} upon it through
actuators, in service of a performance measure it is trying to
optimize. This framing is deliberately broad --- a robot, a
trading bot, a thermostat, and a human all fit --- and the value
of the framing is precisely that it factorizes the design problem
into four pieces (percepts, actions, environment, objective) that
can be reasoned about independently. The rational-agency tradition
treats agency in terms of \emph{means--end} reasoning toward a
specified objective, and judges agents by how well they
approximate that ideal under their compute and information
budgets.

\paragraph{Software-engineering agents.}
Through the 1990s a separate literature developed in
distributed-systems and software-engineering venues, associated
with Michael Wooldridge, Nick Jennings, and the multi-agent
systems community. Their characterization is the one most
LLM-era practitioners are unwittingly recapitulating. A software
agent, on this account, is a component that exhibits some
combination of: \emph{autonomy} (it acts without continuous human
intervention), \emph{reactivity} (it responds to changes in its
environment), \emph{pro-activeness} (it pursues goals rather than
merely responding), and \emph{social ability} (it communicates
with other agents or humans). The list is meant to be a family
resemblance rather than a hard checklist, and the tradition
explicitly accommodates ``weak'' agents that have only some of
these properties.

\paragraph{Reinforcement learning.}
The RL literature inherits a particularly clean formal
abstraction. An agent is whatever sits on one side of the
agent--environment interface: it receives an observation, emits an
action, and receives a scalar reward and a next observation in
return. The interface is the definition. What is on the agent
side --- a tabular policy, a deep network, a tree search, a
language model with tools --- is irrelevant to the formalism, and
this is exactly what makes the framework portable. The price of
the cleanliness is that the framework presupposes the environment
hands out a reward signal, which is a strong assumption that does
not survive contact with most practical agentic deployments.

\paragraph{Philosophy of action.}
The oldest tradition, going back through Davidson and Anscombe to
Aristotle, treats agency as the capacity for \emph{intentional
action}: behavior caused by mental states (beliefs, desires,
intentions) that the agent itself can in some sense own.
Philosophical agency is the demanding sense, and most of what AI
calls an agent fails its tests, deliberately. The reason to
mention it here is that the word ``agent'' carries this
connotation in ordinary language, and the resulting equivocation
is the source of much of the public confusion around AI agents.
When a marketing page describes an agent as ``pursuing goals,''
the philosopher hears a claim about intentionality; the engineer
means a tight loop with a goal variable.

\medskip
\noindent
Each of these uses is internally consistent, and each names
something real. The disagreement among them is not about the
facts; it is about which dimension to foreground. Despite the
divergence, a rough consensus has emerged in the agentic-AI
literature about which properties together qualify a contemporary
system as an ``agent'' in the working sense the rest of this book
will use. Four properties recur~\citep{staufer2026aiagentindex},
and the reader will find some combination of them in nearly every
serious definition produced in the last few years:

\begin{itemize}[leftmargin=1.5em,itemsep=0.3em]

\item \textbf{Autonomy.} The system operates without
  continuous human input on each step. It selects its next action
  on its own and only consults a human at decision points it
  itself flags --- or, in higher-autonomy regimes, not at all.
  Autonomy is a spectrum, not a binary: the same architecture can
  be deployed as a copilot (every action proposed for approval),
  as a delegated worker (a plan presented once, executed
  autonomously, reported on completion), or as a fully unattended
  service. The spectrum is what the practical literature now
  calls levels of autonomy, by analogy with the SAE levels of
  driving automation.

\item \textbf{Goal-directedness over a long horizon.} The system
  is given an objective --- ``investigate this bug,''
  ``refactor this module,'' ``triage today's customer
  tickets'' --- and pursues it across many steps. ``Many'' is
  the operative word: an agent that completes its task in a
  single LLM call is, at best, a one-shot tool. What
  distinguishes an agent is that it strings together steps,
  observes the results of earlier steps, and lets those
  observations shape the next ones. The long horizon is what
  separates an agent from a function call.

\item \textbf{Environment interaction.} The system does not
  merely reason in a vacuum; it reads from and writes to an
  external environment, whether that is a file system, a code
  repository, a web browser, a chat interface, an API zoo, a
  database, or all of the above. The tools of
  Chapter~\ref{ch:toolcalls} are the means by which this
  interaction is implemented; the agent is the controller that
  decides which tool to invoke when.

\item \textbf{Generality.} The system is not a single-purpose
  pipeline. Within its operating envelope it can handle a range
  of tasks it was not explicitly programmed for, by combining the
  primitives at its disposal. This is the property that
  distinguishes an agent from a workflow: a workflow executes a
  predetermined sequence; an agent decides, at each step, what
  the next step should be.

\end{itemize}

\noindent
Throughout this book, when we use the word ``agent'' without
qualification, we mean a system that exhibits all four properties
to a meaningful degree. A chat interface that calls a single
search tool is not an agent in this sense, even though it uses
tool calls. A scripted pipeline that strings together three model
calls is not an agent, even though it has multiple steps. A
shell-using developer assistant that takes a vague request,
inspects a repository, edits files, runs tests, reads the output,
and iterates until done --- that is the prototypical agent of the
mid-2020s, and the architecture that the rest of this chapter
unpacks.

Two caveats are worth stating before moving on. First, the
definition is operational, not metaphysical: it makes no claim
about whether the agent ``really'' has goals in the philosophical
sense, ``really'' understands what it is doing, or has any
inner life worth the name. Those questions are interesting and
unresolved, but the engineer building an agent does not need to
settle them to build well. Second, the four properties are
quantitative rather than binary, and the most useful design work
happens by choosing where on each dimension a given system should
sit. Cranking autonomy up to the maximum is not always the right
choice; broadening generality may sacrifice reliability; long
horizons make every other failure mode worse. The rest of the
agents arc is, in large measure, a discussion of how to make
those trade-offs on purpose.

\section{The Harness}
\label{sec:agents-harness}

The four-property definition of \S\ref{sec:agents-what} tells us
what an agent \emph{is}; it does not tell us what an agent is
\emph{made of}. The architectural question is best approached
through an equation that has been gathering currency in the
practitioner literature, due originally to Vivek Trivedy and
popularized in a synthesis post by Addy
Osmani~\citep{osmani2026harness}:
\begin{equation*}
\textsf{Agent} \;=\; \textsf{Model} \;+\; \textsf{Harness}.
\end{equation*}
The slogan that travels with it --- ``if you're not the model,
you're the harness'' --- does most of the work. A raw language
model is not an agent in any of the senses of
\S\ref{sec:agents-what}: it has no autonomy because someone must
hand it each input; no goal-directedness over a long horizon
because each call is independent; no environment interaction
because text in, text out is the entire interface; and no
generality beyond what fits in a single prompt. The agent
properties are not properties of the model. They are properties
of the \emph{system} the model sits inside. Naming that system,
giving it a noun, makes it possible to talk about it as an
engineering object.

The reader with a background in reinforcement learning will
recognize the shape of the picture immediately, and the analogy
is worth drawing carefully because it both clarifies the harness
concept and exposes a place where the analogy breaks.
\S\ref{sec:agents-what} flagged the RL framing: an agent emits an
\emph{action} into an environment, observes the next state, and
loops. In RL the agent box contains the policy (which maps
observations to actions), and may additionally contain value
functions, world models, and exploration logic; the environment
box contains everything outside that interface. The harness is
the agent box \emph{minus the policy}. The model is the policy;
the harness is the rest of the agent-side machinery --- the
piece that decides what observation to feed the model, parses
the model's output into a structured action, dispatches that
action to the relevant tool, captures the result, decides what
of it to keep, updates whatever scratchpad or memory the next
call will see, and decides whether the loop is done. RL gives us
a clean vocabulary --- observation, action, policy --- and we
will use it. What it elides is that in our setting the harness
also partly constitutes the \emph{environment} the model
perceives: the harness chooses which tools the model is told
about, mounts the filesystem the model reads from, runs the
sandbox the model's commands execute in, and formats the
observation the model gets back. The agent--environment boundary
is sharp on a diagram and porous in practice, because the
harness sits on both sides of it.

\subsection{What is in a Harness}
\label{sec:agents-harness-what}

Listing the components makes the abstraction concrete. The
following inventory follows Osmani's synthesis and is broadly
the same set of moving parts the reader will find described
under different names in writing from Anthropic, the Claude Code
team, the Cursor and Codex teams, and the HumanLayer and
12-factor-agents communities. A harness includes, at minimum:

\begin{itemize}[leftmargin=1.5em,itemsep=0.3em]

\item \textbf{Prompts and instructions.} The system prompt, plus
  any project-level instruction files (\code{CLAUDE.md},
  \code{AGENTS.md}, skill definitions, subagent briefs) that the
  harness prepends to the model's input on every turn. This is
  the most-edited and highest-leverage surface of the entire
  system; the prompting-craft material of
  Chapter~\ref{ch:working} applies here.

\item \textbf{Tools and skills.} The set of callable functions
  the model is told about --- shell, filesystem reads and
  writes, web fetch, code execution, MCP servers, custom
  domain-specific actions --- together with their descriptions.
  Chapter~\ref{ch:toolcalls} treated tool calls as a calling
  convention; from the agent's perspective, the harness is what
  decides which tools exist.

\item \textbf{State and memory.} A scratch surface that survives
  across turns. In a typical coding agent this is the filesystem
  itself, often versioned with git so that experiments can be
  branched and mistakes rolled back. Higher-level memory
  abstractions --- compacted transcripts, knowledge files,
  vector stores --- live here too. Chapter~\ref{ch:context} is
  the relevant deep dive.

\item \textbf{Sandboxes and execution infrastructure.} The
  isolated environment in which the model's tool calls actually
  run: containers, gVisor or Firejail boxes, ephemeral worktrees,
  headless browsers. The sandbox is what makes giving an agent
  shell access tractable rather than catastrophic.
  Chapter~\ref{ch:reliable} treats the security model.

\item \textbf{The control loop.} The driver that calls the
  model, parses out tool invocations, dispatches them, packages
  results back into the model's input, and decides when to stop.
  This is the heart of the runtime and is unpacked in
  \S\ref{sec:agents-harness-loop} below.

\item \textbf{Hooks and middleware.} Deterministic checks that
  fire at specific points in the loop --- before a tool call,
  after a file edit, before a commit. Hooks block destructive
  commands, auto-format generated code, run a typechecker, and
  inject the failures back into the model's next input. The
  harness pattern here is that success is silent and failure is
  verbose: the model only hears about a hook when it failed.

\item \textbf{Orchestration.} Logic for spawning subagents,
  handing off between them, routing across models, and merging
  their work back together. A single-agent harness can omit this
  layer; a multi-agent harness lives here.

\item \textbf{Observability.} Logs, traces, token accounting,
  cost and latency metering, replay tooling. The harness is also
  what makes the agent's behavior debuggable after the fact.

\end{itemize}

\noindent
Two observations follow from the inventory. First, the surface
area is enormous, and almost none of it is the model. Open-source
coding agents --- Claude Code, Cursor, Codex, Aider, Cline ---
can run on the same underlying model and produce strikingly
different behavior, because everything in the list above is theirs
to design. Second, with the inventory in front of us the
``observation--action'' framing of RL becomes operational. The
observation is whatever text the harness has assembled and is
about to send to the model: the system prompt plus instructions,
the relevant slice of state, the previous turn's tool results,
the conversation so far. The action is whatever the harness
parses out of the model's response: a tool call with structured
arguments, a final answer, an exit signal. Everything else in the
list is the machinery that turns the first into the second and
back again.

\subsection{The Control Loop}
\label{sec:agents-harness-loop}

The runtime pattern that ties the inventory together is the same
across nearly every modern agent harness and is structurally
identical to the ReAct loop introduced
by~\citet{yao2023react}:

\begin{enumerate}[label=(\arabic*),leftmargin=2em,itemsep=0.2em]
\item Assemble the next \emph{observation} --- system prompt,
  instructions, current state summary, tool catalog, conversation
  so far, last tool result.
\item Call the model.
\item Parse the model's output. If it contains a tool call,
  extract the tool name and arguments; if it contains a final
  answer, return.
\item Dispatch the tool call through the sandbox, gated by
  pre-call hooks.
\item Capture the result, optionally truncate or offload large
  outputs to the filesystem, gate it through post-call hooks,
  format it as the next turn's observation.
\item Update memory, scratchpads, and any cross-turn state.
\item Check termination conditions (final answer, step budget,
  cost budget, fatal hook failure). If none fired, go to step~1.
\end{enumerate}

\noindent
Two things about this loop are worth flagging because they are
sources of recurring confusion. The first is that it is not the
same as the model's autoregressive loop. The model's loop is
inside step~2 and operates on tokens; the agent's loop is the
outer one and operates on tool calls. The two are nested, not
identical. A model with a 200,000-token context window may take
hundreds of internal token-by-token steps to produce a single
agent step. Conversely, an agent may run for thousands of agent
steps, each of which is one model call, each of which is many
token steps. The dimensions matter independently, and bugs at
each level look very different. The second is that
\emph{step~5 --- formatting the observation} is where most of the
``intelligence'' of a good harness actually lives. Whether a
2,000-line log gets inlined verbatim, truncated to its head and
tail, summarized, or written to disk and replaced by a one-line
pointer is a harness decision, not a model decision. The reader
will find that two harnesses running the same model on the same
task can produce wildly different transcripts because their
step-5 policies differ, and that the gap is sometimes larger
than the gap between two models in the same harness. The variant
in which the model emits \emph{executable code} rather than
JSON-shaped tool calls --- CodeAct~\citep{wang2024codeact} ---
is a step-4-and-5 choice and changes nothing about the overall
loop.

\subsection{The Ratchet}
\label{sec:agents-harness-ratchet}

The harness inventory and control loop together describe the
\emph{structure} of a harness; they do not yet describe how it is
built or evolved. The discipline behind the structure has come to
be called \emph{harness engineering}, and its central habit is a
ratchet: every observed failure becomes a permanent constraint
written into the harness, so that the same failure cannot recur.
Osmani~\citep{osmani2026harness} states the rule starkly: ``every
line in a good system prompt should trace back to a specific,
historical failure.''

The mechanism is concrete. If the agent commits with a
commented-out test, the project's \code{AGENTS.md} gains a line
forbidding it; a pre-commit hook starts flagging \code{.skip(} in
the diff; if there is a reviewer subagent, it gets a rule that
blocks the pattern. If the agent ran a destructive shell command,
a hook blocks the pattern by name on the next attempt and the
sandbox's allowlist is tightened. If the agent got lost in a
forty-step task, the harness gains a planner subagent that splits
work into checkpoints. If the agent stopped early on a still-failing
test, a hook intercepts the stop attempt and injects the failing
output back into the loop. Each addition removes a failure mode
once and for all; each is small; the cumulative effect is what
turns a brittle prototype into a system that ships.

The corollary is that there is no universal best harness. The
correct harness for a particular codebase, team, and task is the
one whose constraints map onto that team's failure history.
Imported templates are useful as starting points but should not
be cargo-culted: a constraint that has not been earned by an
observed failure is paying ongoing cost (in context window, in
maintenance, in opacity) for no benefit. The corresponding
release valve is to delete constraints when a model improvement
renders them redundant, which keeps the harness from accreting
into a fossil record of failures the current model no longer
makes. The result is a continuously edited artifact rather than
a configuration file; harnesses, as Osmani puts it, ``don't
shrink, they move'' --- as models improve, the floor rises but
the ceiling rises faster, and new failure modes appear at the
new horizon for the harness to address.

The remaining sections of this chapter take the components named
here and walk through them with the depth they each deserve.

\section{Levels of Autonomy}
\label{sec:agents-autonomy}

\todo{Section stub. Develop the autonomy spectrum sketched
above into a usable taxonomy: operator, collaborator, consultant,
approver, observer (Feng et al.\ framing, also used by the AI
Agent Index). Walk each level through a worked example
(refactoring a module) so the reader sees what changes at the
interface and what changes in the failure modes.}

\section{The Agent Landscape}
\label{sec:agents-landscape}

\todo{Section stub. Survey what has actually been deployed,
using the AI Agent Index~\citep{staufer2026aiagentindex} as the
empirical spine. Three deployment forms (chat with agentic
tools, browser agents, enterprise agent builders); concentration
on three foundation-model families; the transparency findings
(135/240 safety fields blank, only four agents with
agent-specific safety evaluations). Reproduce or reference the
Index's growth figure to anchor the discussion. Use this section
to motivate the safety and evaluation chapter
(Chapter~\ref{ch:reliable}).}

\section{Why Agents Fail}
\label{sec:agents-fail}

\todo{Section stub. Catalogue of failure modes specific to the
agent loop, distinct from the model-level failures of
\S\ref{sec:mental-failures}: compounding error over long
horizons, context-window saturation, tool-call hallucination,
goal drift, premature termination, and infinite loops.
Forward-reference Chapter~\ref{ch:reliable} for the
operational mitigations.}
